\documentclass[aspectratio=169, 12pt]{beamer}

% Page geometry matching PPTX Widescreen (33.86cm x 19.05cm)
\geometry{paperwidth=33.86cm, paperheight=19.05cm}
\paperwidth=33.86cm
\paperheight=19.05cm
\pdfpagewidth=\paperwidth
\pdfpageheight=\paperheight

\usepackage{icesibeamer}
\usepackage{fontspec}
\usepackage[spanish]{babel}
\usepackage{listings}

% Institutional Font: Plus Jakarta Sans from resources/fonts/
\setmainfont{PlusJakartaSans-Regular.ttf}[
    Path = resources/fonts/,
    BoldFont = PlusJakartaSans-Bold.ttf
]
\setsansfont{PlusJakartaSans-Regular.ttf}[
    Path = resources/fonts/,
    BoldFont = PlusJakartaSans-Bold.ttf
]

% Code listing style
\lstset{
    basicstyle=\ttfamily\fontsize{10pt}{12pt}\selectfont,
    keywordstyle=\color{icesiblue}\bfseries,
    stringstyle=\color{icesiorange},
    commentstyle=\color{icesigreen},
    backgroundcolor=\color{icesigray2!30},
    frame=single,
    rulecolor=\color{icesigray2},
    breaklines=true,
    showstringspaces=false
}

\begin{document}

% ============================================================
% SLIDE 1: Portada
% ============================================================
\titleSlideF{Node.js: Servidores\\con JavaScript}{%
    Entorno de ejecución · Arquitectura Asíncrona · Ecosistema NPM}

% ============================================================
% BLOQUE 1: INTRODUCCIÓN (Slides 2–5)
% ============================================================

% SLIDE 2: ¿Qué es Node.js?
\slideStandard{¿Qué es Node.js?}{%
    Node.js es un \textbf{entorno de ejecución} para JavaScript construido sobre el motor \textbf{V8 de Chrome},
    diseñado para ejecutar código JavaScript fuera del navegador, directamente en el servidor.
    \vspace{0.5cm}
    \begin{itemize}
        \item \textbf{Creador:} Ryan Dahl lo presentó en la JSConf EU de \textbf{2009}.
        \item \textbf{Licencia:} Software libre bajo licencia MIT — código abierto y de uso irrestricto.
        \item \textbf{Motor V8:} El mismo motor de alto rendimiento que impulsa el navegador Google Chrome.
        \item \textbf{Asíncrono por diseño:} Maneja miles de operaciones concurrentes sin crear múltiples hilos de sistema.
        \item \textbf{Multiplataforma:} Disponible para Windows, Linux y macOS con comportamiento consistente.
    \end{itemize}
}

% SLIDE 3: Historia y contexto
\slideSidebarLeftPurple{Historia y Evolución}{%
    \textbf{2009} — Ryan Dahl presenta Node.js en la JSConf EU.
    Inspirado por la frustración con los servidores bloqueantes.
    \vspace{0.4cm}

    \textbf{2011} — Lanzamiento de \texttt{npm} (Node Package Manager),
    transformando el ecosistema JavaScript.
    \vspace{0.4cm}

    \textbf{2015} — Fundación de la Node.js Foundation.
    Lanzamiento de Node.js v4.0 con soporte ES6.
    \vspace{0.4cm}

    \textbf{2018} — Fusión con JS Foundation para crear la
    \textbf{OpenJS Foundation}, gobernanza neutral.
    \vspace{0.4cm}

    \textbf{Hoy (v22+)} — Soporte nativo de ES Modules, TypeScript
    experimental, y mejor integración con herramientas modernas.
}{}

% SLIDE 4: Características clave
\slideTwoCols{Características Clave de Node.js}{%
    \textbf{Técnicas}
    \vspace{0.3cm}
    \begin{itemize}
        \item \textbf{Single-Threaded}: Un único hilo ejecuta el código JS, delegando I/O a libuv.
        \item \textbf{Non-Blocking I/O}: Las operaciones de entrada/salida no detienen el hilo principal.
        \item \textbf{Event-Driven}: La lógica de respuesta se define como callbacks o promesas.
        \item \textbf{Streams}: Procesamiento de datos en fragmentos, sin cargar todo en memoria.
    \end{itemize}
}{%
    \textbf{Ecosistema}
    \vspace{0.3cm}
    \begin{itemize}
        \item \textbf{NPM Registry}: Más de 2.5 millones de paquetes publicados disponibles.
        \item \textbf{Frameworks}: Express, Fastify, NestJS, Koa para backend web.
        \item \textbf{Full-Stack JS}: Mismo lenguaje en frontend (React/Vue) y backend.
        \item \textbf{Comunidad activa}: Una de las comunidades open-source más grandes del mundo.
    \end{itemize}
}

% SLIDE 5: Divisor de Sección
\sectionSlideEBlue{Instalación y\\Configuración del Entorno}

% ============================================================
% BLOQUE 2: INSTALACIÓN (Slides 6–8)
% ============================================================

% SLIDE 6: Requisitos e instalación
\slideStripeTopLeft{Instalación de Node.js}{%
    \begin{columns}[T]
        \begin{column}{0.50\textwidth}
            \textbf{Pasos de Instalación (Windows/macOS)}
            \begin{enumerate}
                \item Visitar \texttt{https://nodejs.org}
                \item Descargar el instalador \textbf{LTS} (Long-Term Support)
                \item Ejecutar el instalador con opciones por defecto
                \item Verificar en terminal: \texttt{node -{}-version}
            \end{enumerate}
            \vspace{0.3cm}
            \textbf{Instalación en Linux (Ubuntu/Debian)}
            \begin{enumerate}
                \item \texttt{curl -fsSL https://deb.nodesource.com/setup\_lts.x | sudo -E bash -}
                \item \texttt{sudo apt-get install -y nodejs}
            \end{enumerate}
        \end{column}
        \begin{column}{0.44\textwidth}
            \textbf{NVM — Recomendado para Desarrollo}
            \vspace{0.2cm}
            \begin{itemize}
                \item \textbf{nvm} (Node Version Manager) permite instalar y cambiar entre múltiples versiones de Node.js fácilmente.
                \item Instalar nvm: \texttt{curl -o- https://raw.githubusercontent.com/nvm-sh/nvm/v0.39.7/install.sh | bash}
                \item Usar: \texttt{nvm install --lts} y \texttt{nvm use 22}
            \end{itemize}
        \end{column}
    \end{columns}
}

% SLIDE 7: Verificación y primeros comandos
\slideStandard{Verificación del Entorno}{%
    \begin{columns}[T]
        \begin{column}{0.72\textwidth}
    Una vez instalado Node.js, valida el entorno con estos comandos en la terminal:
    \vspace{0.3cm}
    \begin{itemize}
        \item \texttt{node -{}-version} — Versión instalada (ej. \texttt{v22.18.0})
        \item \texttt{npm -{}-version} — Versión del gestor NPM
        \item \texttt{node -e "console.log('Hola Node.js')"} — Ejecuta JS directo en terminal
        \item \texttt{node} — REPL interactivo para probar expresiones
    \end{itemize}
    \vspace{0.4cm}
    \textbf{Primer archivo JavaScript en el servidor:}
    \vspace{0.2cm}
    \begin{itemize}
        \item Crear archivo \texttt{app.js} con: \texttt{console.log("Servidor iniciado");}
        \item Ejecutar con: \texttt{node app.js}
    \end{itemize}
        \end{column}
        \begin{column}{0.24\textwidth}
            \vspace{0.8cm}
            \includegraphics[width=\textwidth]{slides/nodejs/assets/icon_terminal.png}
        \end{column}
    \end{columns}
}

% SLIDE 8: Estructura de un proyecto Node
\slideThreeCols{Estructura de un Proyecto Node.js}{%
    \textbf{Archivos Esenciales}
    \vspace{0.3cm}
    \begin{itemize}
        \item \texttt{package.json} — Manifiesto del proyecto: nombre, versión, dependencias y scripts.
        \item \texttt{.env} — Variables de entorno (claves, puertos, URLs de BD).
        \item \texttt{.gitignore} — Excluir \texttt{node\_modules/} del repositorio.
    \end{itemize}
}{%
    \textbf{Directorios Típicos}
    \vspace{0.3cm}
    \begin{itemize}
        \item \texttt{src/} — Código fuente principal de la aplicación.
        \item \texttt{routes/} — Definición de rutas y controladores HTTP.
        \item \texttt{models/} — Esquemas de base de datos o entidades.
        \item \texttt{node\_modules/} — Dependencias instaladas (NO en git).
    \end{itemize}
}{%
    \textbf{Inicializar Proyecto}
    \vspace{0.3cm}
    \begin{itemize}
        \item \texttt{npm init -y} — Crea \texttt{package.json} con valores por defecto.
        \item \texttt{npm install express} — Instala una dependencia.
        \item \texttt{npm install -D nodemon} — Dependencia de desarrollo.
        \item \texttt{npm run dev} — Ejecuta el script \texttt{dev} definido en \texttt{package.json}.
    \end{itemize}
}

% ============================================================
% BLOQUE 3: ARQUITECTURA Y EVENT LOOP (Slides 9–13)
% ============================================================

% SLIDE 9: Divisor de Sección
\sectionSlideEOrange{Arquitectura y\\el Event Loop}

% SLIDE 10: Arquitectura de capas
\slideSidebarLeftBlue{Arquitectura Interna de Node.js}{%
    Node.js está construido sobre una pila de tecnologías bien definidas:
    \vspace{0.4cm}

    \textbf{Motor V8 (Google)} — Compila y ejecuta JavaScript con optimizaciones JIT,
    convirtiendo el código JS en código máquina nativo de alta velocidad.
    \vspace{0.3cm}

    \textbf{libuv} — Biblioteca C de E/S asíncrona multiplataforma.
    Gestiona el Event Loop, los hilos del thread pool y las operaciones de red y disco.
    \vspace{0.3cm}

    \textbf{APIs de Node.js} — Capa de JavaScript que expone módulos nativos como
    \texttt{http}, \texttt{fs}, \texttt{path}, \texttt{events}, \texttt{crypto}, \texttt{stream}.
    \vspace{0.3cm}

    \textbf{Tu aplicación} — El código JavaScript que escribes, apoyado por todo lo anterior.
    \vspace{0.2cm}
    \includegraphics[height=1.2cm]{slides/nodejs/assets/icon_server.png}
}{slides/nodejs/assets/nodejs_architecture.png}

% SLIDE 11: Single Thread vs Multi-Thread
\slideTwoCols{Single-Thread vs Multi-Thread}{%
    \textbf{Servidores Tradicionales (Multi-Thread)}
    \vspace{0.3cm}
    \begin{itemize}
        \item Cada solicitud entrante crea un \textbf{nuevo hilo} de sistema operativo.
        \item Los hilos consumen memoria (~2MB por hilo en Java/PHP).
        \item Con 10.000 conexiones simultáneas: problemas de escalabilidad severos.
        \item Esperas de I/O bloquean el hilo completo (tiempo de CPU desperdiciado).
    \end{itemize}
    \vspace{0.3cm}
    \textit{Ejemplo: Apache HTTP Server, Tomcat (Java)}
}{%
    \textbf{Node.js (Single-Thread + Event Loop)}
    \vspace{0.3cm}
    \begin{itemize}
        \item Un único hilo maneja toda la lógica de JavaScript.
        \item Las operaciones de I/O se delegan a \textbf{libuv} de forma no bloqueante.
        \item Puede manejar \textbf{miles de conexiones concurrentes} con mínima memoria.
        \item El hilo principal siempre está disponible para recibir nuevas solicitudes.
    \end{itemize}
    \vspace{0.3cm}
    \textit{Ideal para: APIs de alta concurrencia, WebSockets, chats en tiempo real}
}

% SLIDE 12: El Event Loop
\slideSidebarLeftBlue{El Event Loop}{%
    \begin{columns}[T]
        \begin{column}{0.65\textwidth}
    El Event Loop es el mecanismo central que hace posible la E/S no bloqueante en Node.js.
    \vspace{0.2cm}

    Funciona ejecutando fases en ciclo continuo:
    \begin{enumerate}
        \item \textbf{Timers:} Callbacks de \texttt{setTimeout} y \texttt{setInterval}.
        \item \textbf{Pending I/O:} Callbacks de operaciones de I/O completadas.
        \item \textbf{Idle/Prepare:} Uso interno de libuv.
        \item \textbf{Poll:} Recupera nuevos eventos de I/O.
        \item \textbf{Check:} Ejecuta \texttt{setImmediate()} callbacks.
        \item \textbf{Close Callbacks:} Limpieza de recursos.
    \end{enumerate}
    \vspace{0.2cm}
    \textit{\texttt{process.nextTick()} y microtareas tienen prioridad máxima entre fases.}
        \end{column}
        \begin{column}{0.3\textwidth}
            \vspace{0.3cm}
            \includegraphics[width=\textwidth]{slides/nodejs/assets/icon_eventloop.png}
        \end{column}
    \end{columns}
}{slides/nodejs/assets/event_loop_diagram.png}

% SLIDE 13: Módulos nativos — Cards
\slideFourCards{Módulos Nativos Esenciales}{%
    \textbf{http / https}
    \newline\newline
    Creación de servidores HTTP/HTTPS y clientes para peticiones externas. Base de cualquier API REST.
}{%
    \textbf{fs (File System)}
    \newline\newline
    Lectura, escritura, creación y eliminación de archivos y directorios. Modo síncrono y asíncrono.
}{%
    \textbf{path / os}
    \newline\newline
    Manipulación de rutas de archivos multiplataforma y acceso a información del sistema operativo.
}{%
    \textbf{events / stream}
    \newline\newline
    EventEmitter para patrones de publicación/suscripción. Streams para procesamiento de datos eficiente en flujo.
}

% ============================================================
% BLOQUE 4: PROGRAMACIÓN ASÍNCRONA (Slides 14–18)
% ============================================================

% SLIDE 14: Divisor de Sección
\sectionSlideEGreen{Programación Asíncrona\\en Node.js}

% SLIDE 15: Evolución del manejo asíncrono
\slideSidebarLeftOrange{Evolución del Código Asíncrono}{%
    JavaScript ha evolucionado su modelo de asincronía en tres generaciones:
    \vspace{0.3cm}

    \textbf{1ª Gen — Callbacks (2009)}
    Funciones que se pasan como argumento para ejecutarse al finalizar una tarea.
    Simple pero propenso al \textit{Callback Hell}: anidación profunda e ilegible.
    \vspace{0.3cm}

    \textbf{2ª Gen — Promises / Thenable (ES6, 2015)}
    Objetos que representan el resultado futuro de una operación. Permiten encadenamiento
    con \texttt{.then()} y manejo de errores unificado con \texttt{.catch()}.
    \vspace{0.3cm}

    \textbf{3ª Gen — Async/Await (ES8, 2017)}
    Syntactic sugar sobre Promises. Escribe código asíncrono con la legibilidad
    del código síncrono. Manejo de errores con \texttt{try/catch} clásico.
}{slides/nodejs/assets/async_evolution.png}

% SLIDE 16: Callbacks — detalle y problema
\slideTwoCols{Callbacks: La Primera Generación}{%
    \textbf{Concepto}
    \vspace{0.2cm}

    Un callback es una función pasada como argumento a otra función, que será invocada
    cuando la operación asíncrona concluya.
    \vspace{0.3cm}

    \textbf{Convención Error-First}
    \begin{itemize}
        \item Primer argumento: el \textbf{error} (o \texttt{null} si no hay error).
        \item Segundo argumento: el \textbf{resultado} de la operación exitosa.
    \end{itemize}
    \vspace{0.3cm}

    \textbf{Patrón}:
    \texttt{fs.readFile('file.txt', (err, data) => \{...\})}
}{%
    \begin{columns}[T]
        \begin{column}{0.58\textwidth}
    \textbf{El Problema: Callback Hell}
    \vspace{0.2cm}

    La anidación de múltiples callbacks genera una pirámide de código inmantenible:
    \vspace{0.2cm}

    \begin{itemize}
        \item Dificulta la lectura y depuración.
        \item Manejo de errores disperso en cada nivel.
        \item Viola el principio de \textit{Separation of Concerns}.
        \item Conocido como \textit{"Christmas Tree"} por la forma del código.
    \end{itemize}
    \vspace{0.2cm}
    \textit{Solución: Promises y Async/Await.}
        \end{column}
        \begin{column}{0.38\textwidth}
            \vspace{0.3cm}
            \includegraphics[width=\textwidth]{slides/nodejs/assets/icon_callback.png}
        \end{column}
    \end{columns}
}

% SLIDE 17: Promises vs Async/Await
\slideTwoCols{Promises y Async/Await}{%
    \textbf{Promises (ES6)}
    \vspace{0.3cm}
    \begin{itemize}
        \item Representa una operación que \textbf{puede completarse o fallar} en el futuro.
        \item Estados: \texttt{pending}, \texttt{fulfilled}, \texttt{rejected}.
        \item Encadenamiento: \texttt{fetch(url).then(r => r.json()).then(data => ...).catch(err => ...)}
        \item \texttt{Promise.all([])} — ejecuta múltiples promesas en paralelo y espera todas.
        \item \texttt{Promise.race([])} — resuelve con la primera promesa que complete.
    \end{itemize}
}{%
    \textbf{Async/Await (ES8)}
    \vspace{0.3cm}
    \begin{itemize}
        \item Palabras clave \texttt{async} y \texttt{await} que simplifican el uso de Promises.
        \item Una función marcada \texttt{async} siempre retorna una Promise implícitamente.
        \item \texttt{await} pausa la ejecución de la función \texttt{async} hasta que la Promise resuelva.
        \item Manejo de errores con \texttt{try/catch} al igual que código síncrono.
        \item Mucho más legible y mantenible que el encadenamiento de \texttt{.then()}.
    \end{itemize}
}

% SLIDE 18: Servidor Web básico
\slideStripeTopRight{Crear un Servidor HTTP en Node.js}{%
    \begin{columns}[T]
        \begin{column}{0.48\textwidth}
            \textbf{Pasos para un servidor básico}
            \begin{enumerate}
                \item \textbf{Importar} el módulo nativo \texttt{http} con \texttt{require('http')}.
                \item \textbf{Crear} el servidor con \texttt{http.createServer(callback)}.
                \item El callback recibe los objetos \texttt{req} (IncomingMessage) y \texttt{res} (ServerResponse).
                \item \textbf{Escribir} la respuesta: estado 200, headers y cuerpo (\texttt{res.end()}).
                \item \textbf{Escuchar} en un puerto: \texttt{server.listen(3000, () => {...})}.
            \end{enumerate}
        \end{column}
        \begin{column}{0.48\textwidth}
            \textbf{¿Qué ocurre en cada petición?}
            \begin{itemize}
                \item El sistema operativo recibe la conexión TCP y notifica a \texttt{libuv}.
                \item \texttt{libuv} coloca el evento en la cola del \textbf{Poll Phase} del Event Loop.
                \item Node.js ejecuta el callback con los datos de la petición.
                \item La respuesta se envía de vuelta al cliente por la misma conexión.
                \item El hilo principal queda inmediatamente libre para otra petición.
            \end{itemize}
        \end{column}
    \end{columns}
}

% ============================================================
% BLOQUE 5: ECOSISTEMA NPM (Slides 19–21)
% ============================================================

% SLIDE 19: Divisor de Sección
\sectionSlideC{El Ecosistema NPM}

% SLIDE 20: NPM en profundidad
\slideStandard{NPM: Node Package Manager}{%
    \begin{columns}[T]
        \begin{column}{0.74\textwidth}
    NPM es el gestor de paquetes oficial de Node.js y el \textbf{registro de software más grande del mundo},
    con más de 2.5 millones de paquetes publicados por la comunidad.
    \vspace{0.3cm}
    \begin{itemize}
        \item \textbf{Instalación:} \texttt{npm install <paquete>} — agrega a \texttt{node\_modules/} y registra en \texttt{package.json}.
        \item \textbf{Dependencias de desarrollo:} \texttt{npm install -D <paquete>} para herramientas del entorno de desarrollo.
        \item \textbf{Scripts personalizados:} Define en \texttt{"scripts"} atajos como \texttt{npm run dev}, \texttt{npm run build}, \texttt{npm test}.
        \item \textbf{package-lock.json:} Congela las versiones exactas de dependencias para builds reproducibles.
        \item \textbf{Seguridad:} \texttt{npm audit} busca vulnerabilidades en el registro CVE.
    \end{itemize}
        \end{column}
        \begin{column}{0.22\textwidth}
            \vspace{0.6cm}
            \includegraphics[width=\textwidth]{slides/nodejs/assets/icon_npm.png}
        \end{column}
    \end{columns}
}

% SLIDE 21: Librerías populares
\slideFourCards{Librerías Clave del Ecosistema}{%
    \textbf{Express.js}
    \newline\newline
    Framework minimalista y flexible para construir APIs REST y servidores web. El más usado del ecosistema Node.js.
}{%
    \textbf{Fastify / NestJS}
    \newline\newline
    Fastify: ultra-rápido con validación integrada. NestJS: framework empresarial con decoradores e inyección de dependencias.
}{%
    \textbf{Mongoose / Prisma}
    \newline\newline
    Mongoose: ODM para MongoDB. Prisma: ORM moderno con tipado fuerte para PostgreSQL, MySQL, SQLite y más.
}{%
    \textbf{Jest / Vitest}
    \newline\newline
    Frameworks de testing para pruebas unitarias e integración. Soporte nativo de mocks, snapshots y cobertura de código.
}

% ============================================================
% BLOQUE 6: COMPARATIVAS Y CASOS DE USO (Slides 22–24)
% ============================================================

% SLIDE 22: Divisor de Sección
\sectionSlideEYellow{Node.js vs\\Otras Tecnologías}

% SLIDE 23: Comparativa con Python y Java
\slideSidebarLeftBlue{Node.js vs Python vs Java}{%
    La elección del entorno de ejecución depende del caso de uso:
    \vspace{0.3cm}

    \textbf{Node.js brilla en:}
    \begin{itemize}
        \item APIs REST de alta concurrencia y bajo tiempo de respuesta.
        \item Aplicaciones de tiempo real: chats, notificaciones, dashboards live.
        \item Microservicios y arquitecturas serverless (AWS Lambda, Vercel).
        \item Herramientas CLI y scripts de automatización de DevOps.
    \end{itemize}
    \vspace{0.3cm}

    \textbf{Python es superior en:} Machine Learning, análisis de datos, scripts científicos.

    \textbf{Java es superior en:} Aplicaciones empresariales, sistemas de misión crítica,
    procesamiento intensivo de CPU, ecosistemas Android/Spring.
}{slides/nodejs/assets/node_vs_others.png}

% SLIDE 24: Cuándo usar / no usar Node.js
\slideTwoCols{¿Cuándo elegir Node.js?}{%
    \textbf{\textcolor{icesigreen}{[SI]}} \textbf{Casos de Uso Ideales}
    \vspace{0.3cm}
    \begin{itemize}
        \item \textbf{APIs REST / GraphQL}: Alta concurrencia, I/O intensivo.
        \item \textbf{Tiempo real}: WebSockets, Socket.io, chats, colaboración.
        \item \textbf{Microservicios}: Arranque rápido, bajo consumo de memoria.
        \item \textbf{BFF}: Backend for Frontend en arquitecturas SPA.
        \item \textbf{Streaming}: Procesamiento de video/audio con Streams.
        \item \textbf{CLI Tools}: npm, webpack, eslint están escritos en Node.
        \item \textbf{Serverless}: Funciones AWS Lambda, Google Cloud Functions.
    \end{itemize}
}{%
    \textbf{\textcolor{icesiorange}{[NO]}} \textbf{Cuándo Evitar Node.js}
    \vspace{0.3cm}
    \begin{itemize}
        \item \textbf{CPU-Bound}: Cálculos matemáticos intensivos bloquean el Event Loop.
        \item \textbf{Machine Learning}: Usar Python (TensorFlow, PyTorch) o servicios cloud.
        \item \textbf{Procesamiento de video/imagen}: FFmpeg + Python/C++ son más adecuados.
        \item \textbf{Big Data}: Ecosistema Spark/Hadoop en Java/Scala es más maduro.
        \item \textbf{Sistemas embebidos}: C/C++/Rust tienen mejor acceso a hardware.
    \end{itemize}
    \vspace{0.3cm}
    \textit{Regla de oro: si la operación es I/O-bound → Node.js. Si es CPU-bound → buscar alternativa.}
}

% ============================================================
% CIERRE (Slides 25)
% ============================================================

% SLIDE 25: Recursos y cierre
\titleSlideA{¿Preguntas?}{%
    nodejs.org · npmjs.com · nodejs.dev/learn}

\end{document}
